% Partie : théorie de la démonstration
% Chapitre : calcul des séquents
% Section : introduction

\documentclass[../../../include/open-logic-section]{subfiles}

\begin{document}

\olfileid[fr]{pt}{seq}{int}

\olsection{Introduction}

Les calculs des séquents sont des systèmes de preuve
introduits par Gerhard Gentzen dans les années~1930.
Gentzen les employait surtout pour étudier la structure
possible des preuves et leurs propriétés, plutôt que
comme méthode pratique de démonstration. Ces questions
sont précisément celles de la théorie de la démonstration.

Ces calculs portent sur des \emph{séquents}, plutôt
que sur des formules isolées. Un séquent est un couple
de suites ou de multiensembles de formules.
Les systèmes originaux de Gentzen, \Log{LK} et \Log{LJ},
utilisaient des suites ; on emploie aujourd'hui souvent
des multiensembles. Un multiensemble ressemble à un
ensemble, mais un même élément peut y figurer plusieurs
fois. L'ordre n'importe pas : $\{!A, !A, !B\}$ et
$\{!A, !B, !A\}$ représentent le même multiensemble.
Il est cependant différent de $\{!A, !B\}$ et
de $\{!A, !A, !B, !B\}$, car les multiplicités
de $!A$ et de $!B$ diffèrent.

On désigne traditionnellement les multiensembles
de formules par des lettres grecques majuscules :
$\Gamma$, $\Delta$, $\Pi$, $\Lambda$, $\Theta$.
Les formules appartiennent au langage logique
considéré. Plutôt que d'écrire le couple
$\tuple{\Gamma, \Delta}$, on sépare ses deux
membres par un symbole de séquent, ici~$\Sequent$.

\begin{defn}[Séquent]
Un \emph{séquent} est une expression de la forme
\[
\Gamma \Sequent \Delta
\]
où $\Gamma$ et $\Delta$ sont des multiensembles
finis, éventuellement vides, de formules.
On appelle $\Gamma$ l'\emph{antécédent}
et $\Delta$ le \emph{conséquent}.
\end{defn}

Soit $!G$ la conjonction des formules de~$\Gamma$,
la conjonction vide valant~$\ltrue$, et soit
$!D$ la disjonction des formules de~$\Delta$,
la disjonction vide valant~$\lfalse$.
Le séquent $\Gamma \Sequent \Delta$ est vrai
dans une structure~$\Struct{M}$ si et seulement
si $!G \lif !D$ y est vrai. En logique classique,
cela signifie qu'au moins une formule de~$\Gamma$
est fausse ou qu'au moins une formule de~$\Delta$
est vraie.

Si $\Gamma$ est un multiensemble de formules,
on note $\Gamma, !A$, ou $!A, \Gamma$, le
résultat de l'ajout d'une occurrence de $!A$
à~$\Gamma$. Si $\Delta$ est aussi un multiensemble,
$\Gamma, \Delta$ désigne leur union additive :
si $!A$ a multiplicité~$n$ dans $\Gamma$,
c'est-à-dire s'il y figure $n$ fois, et multiplicité~$m$
dans $\Delta$, alors sa multiplicité dans
$\Gamma, \Delta$ est $n+m$.
Un membre vide est généralement laissé en blanc.
Ainsi, $\Sequent !A$ a un antécédent vide
et un conséquent contenant $!A$ avec multiplicité~$1$.

Une preuve est un arbre de séquents. Ses séquents
supérieurs, ou initiaux, sont des axiomes du système,
et chaque autre séquent est la conclusion d'une règle
d'inférence dont les prémisses se trouvent immédiatement
au-dessus. Nous considérons d'abord
deux calculs pour la logique classique :
$\Log{G1c}$ et~$\Log{G3c}$. Leurs axiomes et
règles figurent dans les
\cref{pt:seq:int:tab:G1c,pt:seq:int:tab:G3c}.
\oliflabeldef{fol:seq::chap}{ (Le système original
\Log{LK} de Gentzen est présenté dans le
\olref[fol][seq][]{chap}.)}{}

\subfile{rules-G1c}
\subfile{rules-G3c}

Les différences entre ces systèmes ont des
conséquences sur leurs propriétés de preuve.
Les axiomes d'identité de \Log{G1c} sont
$!A \Sequent !A$, sans formule de contexte.
Ceux de \Log{G3c} ont la forme
$!C, \Gamma \Sequent \Delta, !C$, avec des
contextes $\Gamma$ et~$\Delta$ arbitraires,
mais la formule~$!C$ commune aux deux membres
doit être \emph{atomique}.
\Log{G1c} a deux versions de $\LeftR{\land}$
et de $\RightR{\lor}$, tandis que \Log{G3c}
n'en a qu'une de chaque. \Log{G1c} possède
aussi des règles \emph{structurelles} :
la contraction, \LeftR{\Contraction}
et \RightR{\Contraction}, et l'affaiblissement,
\LeftR{\Weakening} et \RightR{\Weakening}.

Les systèmes \Log{G1c} et \Log{G3c} sont
corrects et complets pour la logique classique :
tout séquent qu'ils prouvent est vrai dans
toute structure, et tout séquent vrai dans
toute structure possède une preuve.
La question de l'existence d'une preuve
peut donc aussi être étudiée par des
méthodes sémantiques, notamment celles
de la théorie des modèles. Les deux calculs
prouvent exactement les mêmes séquents.

La théorie de la démonstration étudie toutefois
aussi \emph{comment} on prouve. Peut-on toujours
éviter une règle ? Peut-on l'ajouter sans
obtenir de nouveaux séquents prouvables ?
Peut-on transformer les preuves d'un système
en preuves d'un autre ? Une réponse positive
se démontre souvent par récurrence sur
la complexité ou sur la structure des preuves.
Elle fournit alors non seulement un résultat,
par exemple que \Log{G3c} prouve tout séquent
prouvé par \Log{G1c}, mais aussi une procédure
qui transforme une preuve de~\Log{G1c}
en une preuve de~\Log{G3c}.

Le \emph{théorème d'élimination des coupures}
est un résultat central. On peut ajouter
la règle de coupure
\[
\Axiom$\Gamma \fCenter \Delta, !A$
\Axiom$!A, \Pi \fCenter \Lambda$
\RightLabel{\Cut}
\BinaryInf$\Gamma, \Pi \fCenter \Delta, \Lambda$
\DisplayProof
\]
sans augmenter l'ensemble des séquents
prouvables. Sa démonstration fournit en outre
une méthode qui transforme toute preuve
utilisant les règles de \Log{G1c}, ou
de \Log{G3c}, avec~\Cut{}, en une preuve
qui n'utilise que les règles du calcul
initial. Cette méthode \emph{élimine}
donc les coupures.

Les résultats classiques du premier ordre présentés
ici concernent les séquents de formules closes sans
égalité logique. Les règles quantifiées emploient des
termes clos et des constantes propres. Le modèle de
termes et la recherche de preuves suivent ces
conventions. Pour des formules à variables libres,
il faut aussi préciser les affectations et leurs
effets sur l'instanciation.

La validité caractérise les séquents prouvables.
Au premier ordre, la recherche peut se poursuivre
%
%
indéfiniment sur un séquent invalide.

\end{document}
